%% =====================================================================
%%  T-14-03  The Small First Yes
%%  -------------------------------------------------------------------
%%  One idea per file. Core is mandatory. Slots in canonical order.
%%  Max 700 words. No layout commands. No filler. New source material
%%  for this entry goes in sources/<BOOK>/T-14-03.tex -- never by
%%  rewriting this file (docs/RULES.md R0.1).
%% =====================================================================
%% @kind: topic
%% @id: T-14-03
%% @title: The Small First Yes
%% @subtitle: A commitment escalates on its own fuel
%% @chapter: c14
%% @order: 30
%% @status: stable
%% @sources: B4
%% @updated: 2026-09-19
%% @allow-rewrite: slot migration to the seven-question format (docs/RULES.md section 2)

\topic{T-14-03}{The Small First Yes}{A commitment escalates on its own fuel}


\begin{core}
A person who has agreed to something small will agree to something much larger in the same direction, at a rate the larger request would not achieve on its own. The first agreement is not a step toward the second; it is what makes the second look like the same decision continued.
\end{core}

\begin{mechanism}
Consistency is treated as a virtue and as evidence of a settled character, so once a position has been taken the pressure to hold it is partly social and partly internal. The useful consequence is that the position begins to generate its own support.

After agreeing, a person assembles reasons for having agreed, notices evidence they would previously have ignored, and adopts a self-description that the agreement implies. Those new reasons survive the removal of the original inducement, which is the whole point: the structure holds itself up after the scaffolding is taken away. A small request is therefore valuable in proportion to how far the self-description it implies extends, not in proportion to what it costs.
\end{mechanism}

\begin{application}
  \item Ask for the smallest commitment that still implies the self-description you want installed.
  \item Let the counterpart state the reason for agreeing. Their reason is more durable than yours and is not withdrawn when you leave.
  \item Escalate gradually and always in the same direction; a lateral move reads as a new request and reopens the decision.
  \item Remove the inducement once the position is taken, so that what remains is supported only by the counterpart's own reasons.
\end{application}

\begin{feedback}
  \item The first request is trivially small and framed as though it settles nothing.
  \item Agreeing to it implies a self-description you would not have offered unprompted.
  \item A larger request follows in the same direction and is presented as continuous with the first.
  \item You find yourself producing reasons for the position that were not present when you took it.
  \item Refusing the second request would require explaining the first.
\end{feedback}

\begin{failure}
The technique produces compliance that is durable but shallow: the counterpart holds the position because they have built reasons for it, and those reasons collapse abruptly if the escalation is ever named. Being caught manufacturing a commitment loses the relationship rather than the single transaction, which is a poor trade where you will meet the person again. Used on yourself in reverse it is legitimate --- a small public commitment to something you actually want to do is one of the few reliable self-management tools there is.
\end{failure}

\begin{countermeasures}
  \item Treat every request as the first one. Ask whether you would agree to it having never agreed to anything before.
  \item Separate the two questions: what did I commit to, and what is now being asked? Small agreements are frequently presented as though they covered more ground than they did.
  \item Watch for reasons you have acquired since agreeing. New justifications appearing after a decision are usually the decision defending itself.
  \item Allow yourself to change your mind on new information. The charge of inconsistency is only a cost if you accept that consistency is itself the virtue being tested.
\end{countermeasures}

\sources{B4}
\seesources
\seealso{T-14-01, T-14-04, T-06-04, T-24-02}
